
\section{The one- and two-periodic boundary examples}\label{sec:periodic-ladder}

\subsection{The AHKLY localisation on the left-module side}
For the algebra and presentation fixed in \Cref{sec:preliminaries}, the cone of \(\sigma_L\) is quasi-isomorphic to its cokernel
\[
 U=\langle e_2,\beta\rangle.
\]
Its endomorphism algebra is \(E=k[\eps]/(\eps^2)\), and \(U\cong E\) as a left \(E\)-module.

\begin{theorem}[One-periodic AHKLY defect]\label{thm:ahkly-left}
In left modules,
\[
 \DeltaB(\lambda)\simeq\Dsg(E)\simeq\vect_k,
\]
with suspension isomorphic to the identity.  The comparison with \(\Defsg(\lambda)\) is an equivalence before idempotent completion.  The old perfect-restriction hypothesis fails because \({}_AS\) has infinite projective dimension.
\end{theorem}

\begin{proof}
The right source module \(S_A\) is projective, so \Cref{thm:kernel-dictionary} applies.  The dual-numbers algebra is self-injective local; every finite module is a direct sum of copies of its regular module and its simple \(k\).  The minimal resolution
\[
 \cdots\xrightarrow{\eps}E\xrightarrow{\eps}E\twoheadrightarrow k
\]
has period one.  Hence the stable category is the split category generated by one object fixed by suspension.  It is idempotent complete, so the comparison is an equivalence before idempotent completion.  The supplemental side-verification script gives the infinite periodic left resolution of \(S\).
\end{proof}

\subsection{The opposite right-module boundary}

This section uses right modules.  Passing to opposite rings returns to the left-module convention, and the kernel algebra below is self-opposite.  The underlying recollement example is Angeleri H\"ugel--Koenig--Liu--Yang's Example~4.13(b) \cite{AHKLY17}; the universal-localisation identification, the kernel-dga computation, and the bounded-defect conclusion below are additional.

Let \(A=kQ/I\), where
\[
 Q:\quad 1\xrightarrow{\alpha}2,
 \qquad \beta:2\to2,
\]
and
\[
 \beta^2=0,
 \qquad
 \alpha\beta=0.
\]
Let \(P_i=e_iA\), and let
\[
 \sigma=\alpha:P_2\longrightarrow P_1
\]
be left multiplication by \(\alpha\).  Put
\[
 C=\Cone(\sigma)=[P_2\xrightarrow{\alpha}P_1].
\]

\begin{proposition}\label{prop:dual-localisation}
The universal localisation of \(A\) at \(\sigma\) is
\[
 \lambda:A\longrightarrow M_2(k),
\]
given by
\[
 e_1\mapsto E_{11},\quad
 e_2\mapsto E_{22},\quad
 \alpha\mapsto E_{12},\quad
 \beta\mapsto0.
\]
It is a homological ring epimorphism.
\end{proposition}

\begin{proof}
Left multiplication by \(E_{12}\) gives an isomorphism
\[
 E_{22}M_2(k)\xrightarrow{\sim}E_{11}M_2(k)
\]
with inverse left multiplication by \(E_{21}\), so the map inverts \(\sigma\).

If \(f:A\to T\) inverts \(\sigma\), there is a unique
\(\gamma\in f(e_2)Tf(e_1)\) such that
\[
 f(\alpha)\gamma=f(e_1),
 \qquad
 \gamma f(\alpha)=f(e_2).
\]
Then
\[
 f(\beta)=f(e_2)f(\beta)
 =\gamma f(\alpha)f(\beta)=0.
\]
The elements \(f(e_1),f(e_2),f(\alpha),\gamma\) satisfy the matrix-unit relations and produce a unique factorisation through \(M_2(k)\).  Thus the map is universal.

As a right \(A\)-module, \(M_2(k)\cong P_1^{\oplus2}\).  In particular it is projective, so
\(\Tor_i^A(M_2(k),M_2(k))=0\) for \(i>0\).
\end{proof}

\begin{proposition}\label{prop:dual-dga}
The dg endomorphism algebra of \(C\) is quasi-isomorphic to the ordinary dual numbers
\[
 E=k[\eps]/(\eps^2)
\]
concentrated in degree zero.  As a left \(E\)-complex, the realisation bimodule is
\[
 C\simeq k\oplus k[1].
\]
\end{proposition}

\begin{proof}
The mapping complex \(\End_A^\bullet(C)\) has
\[
 \End^{-1}(C)=\Hom_A(P_1,P_2)=0,
\]
\[
 \End^0(C)=
 \End_A(P_2)\oplus\End_A(P_1)
 =(ke_2\oplus k\beta)\oplus ke_1,
\]
and
\[
 \End^1(C)=\Hom_A(P_2,P_1)=k\alpha.
\]
For \((ae_2+b\beta,ce_1)\in\End^0(C)\), the differential is
\[
 (ae_2+b\beta,ce_1)\longmapsto(c-a)\alpha.
\]
It is surjective.  Its kernel consists of
\((ae_2+b\beta,ae_1)\), and the class of \((\beta,0)\) generates a square-zero radical.  Thus
\[
 H^0\End_A^\bullet(C)\cong k[\eps]/(\eps^2),
 \qquad H^i=0\ (i\ne0),
\]
which gives the quasi-isomorphism.

The element \(\eps\) acts on \(C\) by the chain endomorphism
\((\beta,0)\).  As a left \(E\)-module, \(P_2\cong E\), while
\(P_1=ke_1\oplus k\alpha\) is semisimple and \(\eps\) acts by zero.  The differential sends \(e_2\mapsto\alpha\) and \(\beta\mapsto0\).  Hence
\[
 C\cong[E\twoheadrightarrow k\alpha]\oplus[0\to ke_1]
 \simeq k[1]\oplus k
\]
in \(D(E)\).
\end{proof}

\begin{theorem}[Dual-numbers boundary theorem]\label{thm:dual-numbers}
For the localisation above,
\[
 \Db_\lambda(E)=\per(E).
\]
Consequently
\[
 \DeltaB(\lambda)=0,
 \qquad
 \Dfd(E)/\per(E)=\Dsg(E)\ne0.
\]
Thus a connective ordinary kernel algebra does not determine the bounded localisation defect through its intrinsic finite-dimensional singularity category.
\end{theorem}

\begin{proof}
By \Cref{prop:dual-dga}, for \(Y\in D(E)\), bounded realisation is equivalent to bounded finite-dimensional cohomology of
\[
 Y\Ltensor_Ek,
\]
because
\(Y\Ltensor_EC\simeq(Y\Ltensor_Ek)\oplus(Y\Ltensor_Ek)[1]\).

Choose a pure-minimal semi-projective resolution \(P\to Y\).  Such resolutions exist over the left perfect ring \(E\) by Christensen--Thompson \cite[Theorem~5.7]{ChristensenThompson}; for complexes of projectives over a left perfect ring, pure-minimality agrees with the usual minimality by \cite[Theorem~3.13 and Corollary~3.14]{ChristensenThompson}.  Minimality forces
\[
 d(P^n)\subseteq(\eps)P^{n+1}
\]
for all \(n\): a unit entry in a differential matrix would split off a contractible rank-one two-term summand.  Tensoring with
\(k=E/(\eps)\) therefore kills all differentials, and
\[
 H^n(P\otimes_Ek)=P^n/(\eps)P^n.
\]
If \(Y\Ltensor_Ek\) has bounded finite-dimensional cohomology, then only finitely many \(P^n/(\eps)P^n\) are nonzero and each is finite-dimensional.  Since every \(P^n\) is free over the artinian local ring \(E\), only finitely many \(P^n\) are nonzero and each has finite rank.  Hence \(P\) is a bounded complex of finitely generated projectives and \(Y\in\per(E)\).

Conversely, perfect dg modules clearly realise boundedly.  Therefore
\(\Db_\lambda(E)=\per(E)\), and \Cref{cor:relative-dg-bounded} gives \(\DeltaB=0\).

The simple module \(k\) has the infinite periodic resolution
\[
 \cdots\xrightarrow{\eps}E\xrightarrow{\eps}E
 \xrightarrow{\eps}E\twoheadrightarrow k\to0.
\]
Thus \(k\in\Dfd(E)\setminus\per(E)\), proving the second quotient is nonzero.
\end{proof}

\begin{corollary}[Refutation of cone-only absorption]\label{cor:false-cone-absorption}
Stable flatness together with absorption of the defining cones does not imply relative singularity invariance.
\end{corollary}

\begin{proof}
Take the relative denominator to be the regular module \(M=A\).  Then
\(\Cone(\sigma)\in\Kb(\operatorname{proj}A)=\per(A)\), so the defining cone is absorbed.  The old cone-only assertion would give
\[
 \Dsg(A)\simeq\Dsg(M_2(k))=0.
\]
But the simple right module \(S_2\) has the infinite periodic resolution
\[
 \cdots\xrightarrow{\beta}P_2\xrightarrow{\beta}P_2
 \xrightarrow{\beta}P_2\twoheadrightarrow S_2\to0,
\]
so its class in \(\Dsg(A)\) is nonzero.  The failure occurs because the relevant bounded extension property is absent; after tensoring with \(M_2(k)\), the maps induced by \(\beta\) vanish and the derived tensor is unbounded.  The compact cone does not control the full bounded kernel.
\end{proof}


\subsection{Boundary with Jin--Yang--Zhou}\label{subsec:jyz-boundary}
\begin{lemma}[Fixed-side non-vacuousness]\label{lem:jyz-boundary}
The left-module application of \Cref{thm:kernel-dictionary} to the AHKLY localisation does not satisfy the downward-extension hypothesis of Jin--Yang--Zhou's Theorem~1.1~\cite{JinYangZhou}.  In AHKLY's original right-module convention, the recollement extends downward because \(i_*(S)=P_1\oplus P_1\) is compact.  Passing to opposites changes the controlling module to \({}_AS\), which is not perfect.
\end{lemma}

\begin{proof}
AHKLY compute the right-module recollement in Example~4.13(b)~\cite{AHKLY17}; their compactness criterion makes the displayed right-module object compact.  In the global left convention the same statement is the right-module recollement over \(A^{\op}\), and the object controlling downward extension is \({}_AS\).  The explicit periodic resolution in the supplemental computation has nonzero syzygies in every degree, so \({}_AS\notin\per A\).  AHKLY's correction in Section~4.4 records Example~4.13(b) as a failure of the missing two-sided finite-projective-dimension condition.  Thus the downward hypothesis fails on precisely the side where the Kernel Dictionary applies.
\end{proof}

Even when both theories apply, their roles differ: Jin--Yang--Zhou identify the quotient in a singularity sequence obtained from a downward recollement, while the Kernel Dictionary identifies the compactly generated kernel.

\subsection{Central localisations and the node}
\begin{lemma}[Central nilpotence]\label{lem:central-nilpotence}
Let \(z\in Z(R)\) act trivially on the cohomology of \(X\in\Db(R\text{-}\modcat)\).  If \(X\) has cohomology in at most \(m\) degrees, then \(z_X^m=0\) in \(D(R)\).
\end{lemma}

\begin{proof}
Induct on \(m\).  For maximal \(b\) with \(H^b(X)\ne0\), use the truncation triangle
\[
 \tau_{\le b-1}X\xrightarrow{i}X\xrightarrow{p}H^b(X)[-b]\longrightarrow.
\]
Since \(pz_X=0\), exactness gives \(a:X\to\tau_{\le b-1}X\) with \(z_X=ia\).  Centrality gives \(az_X=z_{\tau_{\le b-1}X}a\), and induction yields \(z_X^m=0\).
\end{proof}

\begin{proposition}[Central-localisation factorisation]\label{prop:central-localisation}
For a finite central localisation \(R\to R[s_1^{-1},\ldots,s_r^{-1}]\), every map from a bounded kernel object to a perfect object factors through the compact kernel.  Hence the comparison is fully faithful and Karoubi dense.
\end{proposition}

\begin{proof}
A common monomial \(u\) in the \(s_i\) annihilates the finitely generated cohomology of a bounded kernel object \(X\).  By \Cref{lem:central-nilpotence}, a power \(u^m\) is zero on \(X\).  For \(f:X\to P\) with \(P\) perfect, centrality gives \((u^m)_Pf=0\).  The multiplication triangle on \(P\) factors \(f\) through \(\Cone(u^m:P\to P)[-1]\), which is perfect and killed by localisation.
\end{proof}

\begin{theorem}[Two-periodic node defect]\label{thm:node}
For
\[
 R=k[[x,y]]/(xy)\longrightarrow R[1/x]\cong k((x)),
\]
the comparison is an equivalence before idempotent completion and
\[
 \DeltaB(\lambda)\simeq\Defsg(\lambda)\simeq\Dsg(R)
 \simeq\vect_k^{\mathbb Z/2}.
\]
Perfect restriction fails.
\end{theorem}

\begin{proof}
The strict cone is \(T=[R\xrightarrow{x}R]\), not \(R/(x)\).  Its endomorphism dga before taking an opposite has
\[
 E^{-1}=R,\quad E^0=R^2,\quad E^1=R,
\]
\[
 d^{-1}(h)=(xh,xh),\qquad d^0(a,b)=x(a-b).
\]
The comparison is full and Karoubi dense by \Cref{prop:central-localisation}.

We classify maximal Cohen--Macaulay modules directly.  This is the standard split-node classification; compare Greuel--Kn\"orrer~\cite{GreuelKnorrer85} under their algebraically closed characteristic-zero convention.  The argument below is included because the present split calculation is valid over an arbitrary field.

Write \(A=k[[x]]\), \(B=k[[y]]\), and
\[
 R=A\times_kB\subseteq\widetilde R=A\oplus B,
 \qquad
 \mathfrak c=xA\oplus yB.
\]
A finite maximal Cohen--Macaulay module \(X\) is torsion-free.  Inside \(Q(R)\otimes_RX\), put
\(\widetilde X=\widetilde R X\cong A^p\oplus B^q\).  Since
\(\mathfrak c\widetilde R\subseteq R\), one has
\[
 \mathfrak c\widetilde X\subseteq X\subseteq\widetilde X.
\]
Thus \(X\) is the inverse image of
\[
 W=X/\mathfrak c\widetilde X
 \subseteq
 \widetilde X/\mathfrak c\widetilde X
 \cong k^p\oplus k^q.
\]
The two projections of \(W\) are surjective: \(X\) generates \(\widetilde X\) over \(\widetilde R\), and Nakayama's lemma applied on each branch gives the assertion.  With
\[
 K_x=W\cap(k^p\oplus0),
 \qquad
 K_y=W\cap(0\oplus k^q),
\]
the subspace \(W\) induces an isomorphism
\[
 k^p/K_x\xrightarrow{\sim}k^q/K_y
\]
and therefore, after choosing complements and bases,
\[
 W=K_x\oplus K_y\oplus\operatorname{graph}(\id_{k^a}).
\]
Every change of basis in \(\operatorname{GL}_p(k)\times\operatorname{GL}_q(k)\) lifts to one in \(\operatorname{GL}_p(A)\times\operatorname{GL}_q(B)\).  Hence this normal form is induced by an automorphism of \(\widetilde X\), not merely by an abstract vector-space isomorphism, and taking inverse images preserves the \(R\)-module isomorphism problem.  The branch coordinates lift to \(R/(y)\) and \(R/(x)\), while each graph coordinate lifts to \(R\).  Consequently
\[
 X\cong R^a\oplus(R/(y))^b\oplus(R/(x))^c.
\]
Thus the nonfree indecomposables are \(N=R/(y)\) and \(M=R/(x)\).  One has \(\End_R(N)=k[[x]]\), and the endomorphisms factoring through a free module form the ideal \((x)\); hence \(\underline{\End}_R(N)=k\).  Similarly \(\underline{\End}_R(M)=k\), while \(\Hom_R(M,N)=\Hom_R(N,M)=0\).  The projective-cover kernels are \((x)\cong N\) and \((y)\cong M\), so \(\Omega M\cong N\) and \(\Omega N\cong M\).  Row and column reduction reduces every stable morphism to identities and zero maps.  Buchweitz's Theorem~4.4.1~\cite{Buchweitz} identifies the singularity category of this Gorenstein ring with the stable MCM category.  Under this equivalence, suspension swaps \(M\) and \(N\); hence the additive equivalence with \(\vect_k^{\mathbb Z/2}\) respects suspension and its split triangulated structure.

The object \(M\) is killed by localisation and is nonperfect.  Its suspension is \(N\), so the image of \(\DeltaB\) contains both generators; Karoubi density is therefore essential surjectivity before idempotent completion.  Finally \(k((x))\) is not finitely generated over \(R\), so it is not perfect.
\end{proof}